% Target: rmp-iv-intersection-rhs-two
% Formula: Applying the inverse of B leaves R, an exposed site, and C.
% Ink: The two virtual indices exposed by the inverse and C sit above R.
% Boundary: The outer end closes through R; two virtual and two physical legs open.
% Source: paper TN-Review-main.tex line 2059; author source ImagesReview Section
%   IV/ImagesReview Section IV.tex lines 87-94 (Rtensor2)
% Capabilities: free-graph, closure, typed-ports
\[
  \begin{tenkzfree}
    \tnput[box, ports={west:virtual,east:virtual,north:physical}]
      {c}{(8mm,0mm)}{C}
    \tnput[pill, ports={west:virtual,east:virtual,135:physical}]
      {r}{(0mm,-8mm)}{\qquad R\qquad}
    \tnput[boundary, ports={south:physical}]{rp}{(-8mm,9mm)}{}
    \tnput[boundary, ports={south:virtual}]{rv}{(-2mm,9mm)}{}
    \tnput[boundary, ports={south:virtual}]{cv}{(2mm,9mm)}{}
    \tnput[boundary, ports={south:physical}]{cp}{(8mm,9mm)}{}
    \tnjoin[route=arc,out=0,in=0]{c.east}{r.east}
    \tnjoin[route=hv]{r.west}{rv.south}
    \tnjoin[route=hv]{c.west}{cv.south}
    \tnjoin[route=hv]{r.135}{rp.south} \tnjoin{c.north}{cp.south}
  \end{tenkzfree}
\]
